\documentclass[a4paper,11pt]{article}
\usepackage[margin=27mm]{geometry}
\usepackage{fontspec,amsmath,amssymb,amsthm}
\setmainfont{STIXGeneral}
\usepackage{unicode-math}\setmathfont{STIX Two Math}
\setlength{\emergencystretch}{2em}
\usepackage[hidelinks]{hyperref}
\newcommand{\Sym}{\operatorname{Sym}}
\newcommand{\gen}[1]{\langle #1\rangle}
\newtheorem{theorem}{Theorem}\newtheorem{lemma}[theorem]{Lemma}\newtheorem{corollary}[theorem]{Corollary}
\title{Two copies of a tree and their permutation intersections}
\author{Human Compression project}\date{Private proof draft; independent review pending}
\begin{document}\maketitle
\begin{abstract}
Synchronizing transpositions in two copies of a tree gives a diagonal
permutation group. Deleting a connected set from one copy releases independent
symmetric factors exactly in the components meeting that set. We prove this
for every subset of the edge labels, and deduce the intersection property
without a rank induction. Adjoining the interchange of the two copies gives
a second family, with full wreath-product group. Path cases recover the unequal-simplex and two-row wreath-product lemmas of
Leemans--Mulpas. A separate application gives the remaining positive family
and their already known uniform counterexample. The proofs here do not use
finite sample calculations.
\end{abstract}

\section{Partitions and a connected deleted set}
Let $T$ be a finite nonempty tree, with vertex set $V$ and edge set $E$.
For $I\subseteq E$, write $\pi_I$ for the partition of $V$ into the connected
components of the spanning forest $(V,I)$. A permutation preserves a partition
in this note if it preserves \emph{each} block, rather than merely permuting
the blocks. The common refinement of two partitions consists of their nonempty
block intersections.

\begin{lemma}\label{lem:paths}
For any $I,J\subseteq E$,
$\pi_I\wedge\pi_J=\pi_{I\cap J}$.
Suppose $Z$ is a nonempty connected set of vertices of $T$, and put
\[
 R_I=\{v\in V:\text{the $I$-component of $v$ meets $Z$}\}.
\]
Then $R_I\cap R_J=R_{I\cap J}$.
\end{lemma}
\begin{proof}
Two vertices lie in the same $I$-component precisely when every edge of their
unique path in $T$ belongs to $I$. Requiring this both for $I$ and $J$ is the
same as requiring it for $I\cap J$, proving the partition identity.
For the second assertion, every vertex of $Z$ lies in all three sets.
Each component of $T\setminus Z$ attaches to $Z$ by exactly one edge:
there is at least one because $T$ is connected, and two would give a cycle
using a path inside that component and the path inside connected $Z$.
Thus for $v\notin Z$ there is a unique first vertex $z(v)$ of $Z$ met by a
path from $v$ to $Z$. The component of $v$ meets $Z$ in $(V,I)$ precisely
when all edges of the path from $v$ to $z(v)$ belong to $I$.
The same simultaneous-path argument proves the assertion.
\end{proof}

\section{Independent factors created by deletion}
Use two disjoint copies $A=\{a_v:v\in V\}$ and
$B=\{b_v:v\in V\}$. For a fixed nonempty connected $Z\subseteq V$ and
an edge $e=\{u,v\}$, define
\[
 t_e=\begin{cases}
 (a_u\ a_v)(b_u\ b_v),&u,v\notin Z,\\
 (a_u\ a_v),&\{u,v\}\cap Z\ne\varnothing.
 \end{cases}
\]
The points $b_z$ with $z\in Z$ are fixed by every mark and may be deleted
from the permutation domain. Put $H_I=\gen{t_e:e\in I}$.
Let $I^Z$ be obtained from $I$ by removing every edge with an endpoint in $Z$.
Identify permutations of either row with permutations of $V$.

\begin{theorem}\label{thm:deleted}
The group $H_I$ consists exactly of the pairs $(\sigma,\tau)\in\Sym(V)^2$
satisfying
\begin{enumerate}
\item $\sigma$ preserves every block of $\pi_I$;
\item $\tau$ preserves every block of $\pi_{I^Z}$;
\item $\sigma(v)=\tau(v)$ for every $v\notin R_I$.
\end{enumerate}
Equivalently, for an $I$-component $C$ disjoint from $Z$, the two rows of $C$
carry one diagonal copy of $\Sym(C)$. For an $I$-component $C$ meeting $Z$,
the first row carries the independent group $\Sym(A_C)$, and the second
row carries an independent symmetric group on every component of
$(C\setminus Z,I^Z)$; each $b_z$ stays fixed. All these factors for different
components are independent.
Consequently $(t_e)_{e\in E}$ is a C-marking, and its full group is
\[
 \Sym(A)\times\prod_{D\text{ component of }T\setminus Z}\Sym(B_D).
\]
\end{theorem}
\begin{proof}
Different components of $(V,I)$ use disjoint sets of edge labels and disjoint
supports, so their groups form a direct product. If $C$ avoids $Z$, each mark
on $C$ acts by the same transposition in both copies. Edge transpositions of
a connected tree generate its full symmetric group: conjugation along a path
produces all transpositions with one fixed endpoint. Thus this component gives
exactly the stated diagonal copy.

Suppose now that $C$ meets $Z$, and choose $z\in C\cap Z$.
We prove that the supported group $\Sym(A_C)$ lies in $H_I$.
Grow a connected set $W\subseteq C$ from $\{z\}$ along its tree edges.
Suppose the group already contains every transposition $(a_z\ a_x)$ for
$x\in W\setminus\{z\}$, and take an edge $\{x,y\}\in I$ with $x\in W$
and $y\notin W$. If $x=z$, that edge meets $Z$ and its mark is precisely
the pure transposition $(a_z\ a_y)$. Otherwise its mark fixes $a_z$, and
conjugating $(a_z\ a_x)$ by it gives $(a_z\ a_y)$, with no action on the
second row. Induction supplies all these transpositions, which generate
$\Sym(A_C)$.

Cancel from each doubled edge mark its now available first-row transposition.
This supplies every second-row edge transposition on $I^Z$ inside $C$.
They generate the symmetric group of each remaining second-row component.
Conversely each original mark belongs to the product of the factors just
listed. Hence we have both containments and the exact factorization. This
also proves the membership description: on components outside $R_I$ the
restrictions must agree, and on components inside $R_I$ they are exactly the
independent permutations permitted by the two partitions.

To intersect $H_I$ and $H_J$, intersect these three literal conditions.
Preservation of both partitions is preservation of their common refinement.
Lemma~\ref{lem:paths} gives the first-row partition $\pi_{I\cap J}$ and
the second-row partition $\pi_{I^Z\cap J^Z}=\pi_{(I\cap J)^Z}$.
The two agreement conditions together require agreement outside
$R_I\cap R_J=R_{I\cap J}$. These are precisely the conditions for
$H_{I\cap J}$, so every subgroup intersection is correct.
The marks are distinct nonidentity involutions because their first-row
transpositions are distinct. Finally $I=E$ has a single component and meets
$Z$, giving the full-group formula.
\end{proof}

If no vertices are deleted, the same elementary argument gives a diagonal
$\Sym(V)$, with all subgroup intersections governed by the partitions
$\pi_I$. The argument also works for any number of identical copies.
The connectedness hypothesis on a nonempty deleted set has content.
For the path $1$--$2$--$3$--$4$ and $Z=\{1,4\}$, the marks are
\[
 (a_1a_2),\quad(a_2a_3)(b_2b_3),\quad(a_3a_4).
\]
Both subgroups generated by consecutive pairs contain the pure transposition
$(b_2b_3)$: cubing the product of the two marks gives it, since the first-row
product has order three. Their common-label subgroup has order two and its
nonidentity element moves the first row, so it does not contain this pure
second-row transposition. The intersection property fails. This is the
unequal-start simplex counterexample displayed in the source on physical
page~9, written in coordinates, not a new source correction.

\section{Interchanging the rows}
Fix one vertex $z\in V$, apply the preceding construction with $Z=\{z\}$,
and retain both complete rows, including the formerly fixed point $b_z$.
Adjoin the involution
\[
 s=\prod_{v\in V}(a_v\ b_v).
\]
Put $D_I=z^{\pi_I}$, the block of $\pi_I$ containing $z$.

\begin{theorem}\label{thm:wreath}
The marks $\{t_e:e\in E\}\cup\{s\}$ form a C-marking of
$\Sym(V)\wr C_2$ in its two-row imprimitive action.
More precisely, write $K_I$ for the row-preserving pairs that each preserve
$\pi_I$ and agree outside $D_I$. Then
\[
 \gen{H_I,s}=K_I\rtimes\gen{s},\qquad
 K_I=\Sym(A_{D_I})\times\Sym(B_{D_I})
       \times\prod_{C\in\pi_I,\ C\ne D_I}\operatorname{diag}\Sym(C).
\]
\end{theorem}
\begin{proof}
Theorem~\ref{thm:deleted} already gives $\Sym(A_{D_I})$ inside $H_I$.
Conjugating by $s$ gives $\Sym(B_{D_I})$. The other components act diagonally
and are fixed by conjugation with $s$. The displayed product is therefore
contained in the row-preserving part of $\gen{H_I,s}$. Conversely this product
contains $H_I$, is normalized by $s$, and is row preserving, whereas $s$
interchanges the two nonempty rows. This proves the asserted semidirect
product and its exact row-preserving kernel. This includes $D_I=\{z\}$,
when its two symmetric factors are trivial.

The path argument in Lemma~\ref{lem:paths} gives
$D_I\cap D_J=D_{I\cap J}$. The membership conditions in the definition of
$K_I$ therefore give $K_I\cap K_J=K_{I\cap J}$. Moreover,
\[
 H_I\cap K_J=H_{I\cap J}.
\]
Indeed, on the first row the required partition is $\pi_{I\cap J}$.
On the second row it is the common refinement of $\pi_{I^{\{z\}}}$ and
$\pi_J$, which is $\pi_{(I\cap J)^{\{z\}}}$ by the unique-path argument.
The agreement region is again the complement of $D_I\cap D_J$.
These are exactly the three conditions of Theorem~\ref{thm:deleted}.

Two selected subgroups neither of which contains the label $s$ intersect
correctly by that theorem. If just one contains $s$, an element of the other
is row preserving, so the last displayed mixed-intersection equality applies.
If both contain $s$, every element has a unique form $ks^\epsilon$, with
$\epsilon\in\{0,1\}$. Because the same $s$ belongs to both groups, their
intersection is $(K_I\cap K_J)\rtimes\gen{s}$, giving precisely the subgroup
on the common labels. Thus every subset intersection is proved. The new mark
is different from all edge marks because it interchanges the rows. For $I=E$,
we have $D_I=V$, so the whole group is the full stated wreath product.
\end{proof}

\section{Independent flips within paired vertices}
There is a companion construction in which one interchanges a single pair,
rather than interchanging the two entire rows. It is convenient to record
its complete subgroup formula alongside the preceding one.

\begin{theorem}\label{thm:binary}
On the two copies of a tree $T$, put
$d_{\{u,v\}}=(a_u a_v)(b_u b_v)$ for every edge and
$f=(a_zb_z)$ at a fixed vertex $z$. These marks form a C-marking of
$C_2\wr\Sym(V)$, in its action on the pairs $\{a_v,b_v\}$.
For selected edge set $I$, the subgroup omitting $f$ is the diagonal
permutation group preserving $\pi_I$. If $f$ is selected, it is
\[
 (C_2)^{D_I}\rtimes\prod_{C\in\pi_I}\Sym(C),
 \qquad D_I=z^{\pi_I},
\]
where $(C_2)^{D_I}$ consists of all independent flips of the pairs indexed
by $D_I$, and each symmetric factor permutes the corresponding pairs.
\end{theorem}
\begin{proof}
Every generated permutation preserves the set of pairs. Such a permutation
has unique coordinates: a permutation of the pair indices, and the subset
of pairs whose orientations it reverses. The selected edge marks generate
exactly the index-permutation group preserving $\pi_I$, with no flips.
Conjugating $f$ by those permutations gives the individual flip at every
vertex of $D_I$; these commute and generate all subsets of flips there.
Conversely their group is preserved by all selected index permutations,
so no further flip can be produced. This proves both formulas and their
actual permutation realization.

In these unique coordinates, intersecting two selected groups intersects
the index-permutation groups and the permitted flip supports. The former
intersection is controlled by $\pi_I\wedge\pi_J=\pi_{I\cap J}$.
If both groups select $f$, the latter support is
$D_I\cap D_J=D_{I\cap J}$; if either omits it, no flips survive.
These are exactly the groups on the common labels. All marks are distinct
nonidentity involutions, proving the C-marking assertion. Taking every edge
makes $D_I=V$ and gives the stated full wreath product.
\end{proof}

For a path with $z$ at an endpoint, put $f$ first, followed by the path
edges starting at $z$. Distant marks commute, so this gives the usual
string marking of $C_2\wr S_{|V|}$. It is the folded-path marking of
L010 Lemma~4.5 (physical pages~9--10), with no exceptional finite base case.
The binary permutation action and its semidirect-product coordinates are
standard; the explicit all-subset formula is the proof needed here.

\section{A terminal extension of a path}
A first new tail edge needs an additional argument: the row swap is a product
of several transpositions. It is not the pure endpoint mark required by a
pure-leaf gluing theorem. The nesting of the terminal components of a path
supplies precisely the extra control.

\begin{theorem}\label{thm:tail}
Let $T$ be the path on $\{1,\ldots,n\}$, let $z=n$, and take the marking
of Theorem~\ref{thm:wreath}. Add a new point $c$ and the last mark
$t=(b_n\ c)$. Order the path edge marks towards $n$, followed by $s,t$.
The result is a string C-marking of $S_{2n+1}$.
Any further path of new pure transpositions attached successively at the
new endpoint again gives a string C-marking, of the full symmetric group
on the resulting points.
\end{theorem}
\begin{proof}
For $I\subseteq E$, let $Q_I^{\epsilon\eta}$ be the selected subgroup,
where $\epsilon$ records whether $s$ is included and $\eta$ records whether
$t$ is included. We already know
\[
 Q_I^{00}=H_I,\qquad Q_I^{10}=K_I\rtimes\gen{s},\qquad
 Q_I^{01}=H_I\times\gen{t}.
\]
The last formula holds because $H_I$ fixes both $b_n$ and $c$.
Write $D_I$ for the terminal interval of the selected path containing $n$,
and put
\[
 R_I=\{c\}\cup A_{D_I}\cup B_{D_I},\qquad
 u_I=\prod_{v\notin D_I}(a_v\ b_v),\qquad
 \Delta_I=\prod_{C\in\pi_I,\ C\ne D_I}\operatorname{diag}\Sym(C).
\]
Then
\begin{equation}\label{eq:both}
 Q_I^{11}=\Sym(R_I)\times\bigl(\Delta_I\times\gen{u_I}\bigr).
\end{equation}
Indeed, the orbit of $b_n$ under $Q_I^{10}$ is $A_{D_I}\cup B_{D_I}$,
and that group fixes $c$. Conjugating $t$ by it gives all transpositions
$(c\ x)$ on $R_I$. These generate its supported symmetric group.
Cancelling restrictions on $R_I$ shows that the induced complementary action
is an independent factor. The formula for $Q_I^{10}$ identifies this
complementary action as $\Delta_I\times\gen{u_I}$, proving
\eqref{eq:both}. An empty complement gives the trivial factor.

We now check all selected-subgroup intersections. The groups with neither
new mark, and those with only $s$, already have the required intersections
by Theorem~\ref{thm:wreath}. A subgroup containing $t$ but not $s$ is
$H_I\times\gen{t}$; its elements move $c$ precisely when their $t$-exponent
is one. Intersecting two such subgroups therefore gives
$H_{I\cap J}\times\gen{t}$. Intersecting one with any group omitting $t$
forces that exponent to be zero, and the preceding mixed-intersection
formula gives $H_{I\cap J}$ when necessary.

For comparisons involving $Q_I^{11}$, its elements which preserve both rows
and fix $c$ form exactly $K_I$. This follows from \eqref{eq:both}: the root
symmetric group then preserves its two row parts, and the complementary
flip is excluded unless it is already trivial. Hence
\[
 Q_I^{11}\cap H_J=H_{I\cap J},\qquad
 Q_I^{11}\cap Q_J^{10}=Q_{I\cap J}^{10},\qquad
 Q_I^{11}\cap Q_J^{01}=Q_{I\cap J}^{01}.
\]
For the second equality, remove the common element $s$ when the row-exchange
exponent is one, and use $K_I\cap K_J=K_{I\cap J}$.
For the third, remove the common element $t$ when the $c$-moving exponent is
one, and use the first equality. These removals are legitimate because both
$s$ and $t$ belong to every $Q_I^{11}$.

It remains to intersect two groups of the last kind. Terminal intervals of
a path are nested, so assume $D_I\subseteq D_J$. Their common selected
component is then $D_{I\cap J}=D_I$. Let
$g\in Q_I^{11}\cap Q_J^{11}$. Its action outside $R_I$ has a single
row-exchange exponent, by \eqref{eq:both}. Multiplying by the common element
$s$, if necessary, makes that exponent zero. If the complement is empty
there is no exponent to remove. The resulting element preserves $R_I$;
cancel its restriction there using $\Sym(R_I)$, which belongs to
$Q_{I\cap J}^{11}$. Call the remainder $g'$.

Now $g'$ fixes $R_I$ pointwise and acts diagonally outside it, preserving
$\pi_I$. Membership in $Q_J^{11}$ forces preservation of every block of
$\pi_J$ as well: outside $D_J$ this is explicit in \eqref{eq:both}; on
$D_J$ the element preserves its union of row copies, and its already
established row preservation makes it preserve $D_J$ in each row.
Thus its common row permutation preserves
$\pi_I\wedge\pi_J=\pi_{I\cap J}$. It fixes $D_I$ and so lies in
$\Delta_{I\cap J}$. Restoring the cancelled root permutation and the possible
$s$ proves $g\in Q_{I\cap J}^{11}$. The reverse inclusion is automatic.
This settles every pair of subsets.

The string relations hold because $t$ is disjoint from every edge mark
(they all fix $b_n$) and is consecutive with $s$. Distinctness and
nonidentity of the marks are immediate from their different supports.
For $I=E$, \eqref{eq:both} is the full $S_{2n+1}$.
Finally the point $c$ is moved only by the pure endpoint transposition $t$.
The proved pure-leaf gluing theorem therefore allows one further leaf at
$c$, and then repeated leaves. Its same transposition-conjugation argument
identifies each connected enlarged group as the full symmetric group.
\end{proof}

The terminal-path hypothesis cannot be silently replaced by an arbitrary
rooted tree. For example, take edges $\{1,2\}$ and $\{1,3\}$, root $1$,
and the four analogous marks $(a_1a_2),(a_1a_3),s,(b_1c)$.
The selected groups with labels $\{\{1,2\},s,t\}$ and
$\{\{1,3\},s,t\}$ both contain $(a_3b_3)$: it is the complementary
flip of the first group and lies in the supported root symmetric factor of
the second. Their common-label subgroup $\gen{s,t}$ is the product of
$\Sym\{a_1,b_1,c\}$ with the single outside flip
$(a_2b_2)(a_3b_3)$, so it does not contain $(a_3b_3)$.
The intersection property therefore fails. The factorization used here
follows by the same transposition conjugation as \eqref{eq:both}; only the
intersection step requires nested terminal components.

\section{Exact path applications and remaining scope}
In the source of Leemans--Mulpas, arXiv:2504.17535v1, Lemma~4.4 (physical
page~9) has two aligned simplex paths on $r+1$ and $h+1$ vertices, with
$1\leq h\leq r-1$. Take $T$ to be the path on $r+1$ vertices and
$Z=\{h+2,\ldots,r+1\}$. The surviving second path has exactly the same
labels as the initial segment of the first. Theorem~\ref{thm:deleted} gives
all intersections and the full group $S_{r+1}\times S_{h+1}$; deleting the
fixed second-row points changes no group or marking. Consecutive edge labels
satisfy the string relations, since nonconsecutive path edges have disjoint
endpoints. The source attributes this lemma to Fernandes--Leemans--Mixer,
\emph{All alternating groups $A_n$ with $n\geq12$ have polytopes of rank
$\lfloor(n-1)/2\rfloor$}, SIAM J. Discrete Math. 26 (2012),
Lemma~6.1;
our argument is a complete local proof of that retained conclusion, not a
claim of priority. Identical paths give its Lemmas~4.2--4.3 by the diagonal
case above.

For Lemma~4.9 (physical pages~12--14), let $T$ be the path on $r$ vertices,
$z=r$, order its edge marks from $1$ towards $r$, and put the row swap last.
The first $r-2$ edge marks are doubled, the last edge mark is a pure
first-row transposition, and $s$ is the full row swap. All nonconsecutive
marks commute: $s$ commutes with every doubled mark and is consecutive with
the sole pure one. Theorem~\ref{thm:wreath} therefore gives the full string
C-marking of $S_r\wr C_2$, including the stated range $r\geq3$.
For the source's $r=3$ coordinates take
$(a_1,a_2,a_3)=(1,2,3)$ and $(b_1,b_2,b_3)=(5,6,4)$.
The marks are exactly
\[
 (12)(56),\qquad(23),\qquad(15)(26)(34).
\]
These are the literal source generators. This uniform argument
replaces its finite base check and induction. No finite test is used as a
premise of our proof.

The graph identification in the last paragraph must always use the literal
row pairs; abstract equality of the wreath-product group alone would not
identify a marking. For Proposition~4.10, put $n=i+1$. Theorem~\ref{thm:wreath} gives its
initial rank-$i+1$ marking on $2i+2$ points. Theorem~\ref{thm:tail}
supplies the first new tail edge; this step is established explicitly because
the initial last mark is the compound row swap. Further pure leaves give
rank $r>i+1$ and degree $r+i+1$, with full symmetric group. This recovers
the complete positive family without using a finite base case.
The auxiliary two-component family, the other positive family and the uniform
counterexample of Section~4 are proved in the separate application. The present
paper gives a shared argument; the names and origins of the recovered source
results remain those just stated. No historical novelty assertion is made.

\section{Formal scope and a character test}
The file \texttt{ParallelIntersections.lean} verifies the literal intersection
of the three permutation constraints in Theorem~\ref{thm:deleted}. It treats
row preservation as an arbitrary binary relation and proves that intersecting
the two conditions refines that relation by conjunction. The agreement regions
combine by intersection. The tree path identities and the exact generation
formulas are proved above; they are not formalized in that file.

The same file verifies a second elementary fact used in the application.
If homomorphisms $\chi,\psi:G\to M$ agree on a generating set $S$, they agree
on $\langle S\rangle$. Thus an element $u$ with
$\chi(u)\ne\psi(u)$ is outside that subgroup. This is the usual equalizer
argument: the equality set is a subgroup, since both maps preserve products
and inverses. Our formal proof explicitly reuses mathlib's theorem that two
homomorphisms agreeing on a set agree on its subgroup closure. The source's
counterexample is obtained when one character is the sign of a local action
and the other records an outside involution. The particular restrictions,
sign values and conjugating elements remain written application proofs.
The earlier six-declaration orbit module is also reused with its original
limited scope; none of these statements alone formalizes a complete gluing
construction.
\end{document}
